\PassOptionsToPackage{table,xcdraw}{xcolor}
\documentclass[11pt, a4paper]{article}

\usepackage[T1]{fontenc}
\usepackage[utf8]{inputenc}
\usepackage{tgpagella}
\usepackage[spanish, es-noshorthands]{babel}
\usepackage{amsmath, amssymb}
\usepackage[most]{tcolorbox}
\usepackage{tabularx}
\usepackage{tikz}
\usepackage[american]{circuitikz}
\usepackage[margin=2cm]{geometry}
\usepackage{fancyhdr}

\pagestyle{fancy}
\fancyhf{}
\setlength{\headheight}{16pt}
\lhead{\textbf{UCB Sede Tarija} | Ing. Mecatrónica}
\rhead{IMT-221 | Diagnóstico Individual EC3 (E1)}
\renewcommand{\headrulewidth}{0.4pt}
\definecolor{ucbblue}{RGB}{0, 51, 102}

\begin{document}

\begin{tcolorbox}[colback=ucbblue!5, colframe=ucbblue, title=\textbf{Hoja Diagnóstica Individual (10 min) — Evidencia de Pérdidas}]
    Acredite su comprensión de las pérdidas por conmutación antes de avanzar a la especificación de componentes físicos para la Etapa 1 (E1)[cite: 16].
\end{tcolorbox}

\noindent\textbf{Nombre Estudiante:} \rule{7cm}{0.15mm} \hfill \textbf{Fecha:} \rule{3cm}{0.15mm} \\[0.3cm]
\textbf{Estado de su evidencia LTspice W12S02 (Marque uno)[cite: 16]:} \\
$\square$ Observación LTspice ejecutada y documentada \quad $\square$ Predicción teórica \quad $\square$ Ninguna corrida

\vspace{0.4cm}
\textbf{1. Físico-Matemática de Transiciones (Turn-ON / Turn-OFF)[cite: 16]} \\
Durante el flanco de encendido de un NMOS en configuración Low-Side con carga resistiva, ¿qué ocurre cualitativamente con $V_{GS}, V_{DS}$ e $I_D$? ¿Qué magnitudes eléctricas "coexisten" peligrosamente durante este breve intervalo transitorio?[cite: 16]
\vspace{2.5cm}

\textbf{2. Definiciones de Potencia y Energía[cite: 16]} \\
Defina matemáticamente las siguientes magnitudes, indicando explícitamente sus unidades físicas en el SI:
\begin{itemize}
    \item Potencia instantánea disipada por el MOSFET ($p_D(t)$)[cite: 16]: \hrulefill
    \item Energía de un evento único de conmutación en la ventana $[t_a, t_b]$ ($E_{event}$)[cite: 16]: \hrulefill
    \item Potencia media de \textit{switching} repetitivo ($P_{sw}$)[cite: 16]: \hrulefill
\end{itemize}

\vspace{0.2cm}
\textbf{3. La Trampa de la "Carga Ideal"[cite: 16]} \\
En una simulación didáctica con $V_{bus} = 12\,\text{V}$ y $R_L = 10\,\Omega$, un cálculo arroja $14.4\,\text{W}$[cite: 16]. 
Identifique a qué variable física corresponde exactamente ese valor y justifique mediante una ley de circuitos. ¿Por qué es un error fatal de diseño asumir que esos $14.4\,\text{W}$ son la disipación térmica del MOSFET real?[cite: 16]
\vspace{3cm}

\textbf{4. Auditoría de Simulación[cite: 16]} \\
Declare una limitación fundamental del modelo pedagógico LTspice (\texttt{MDEMO}) que usamos la semana pasada, y proponga qué medición en el laboratorio pondría a prueba dicha limitación[cite: 16].
\vspace{2.5cm}

\begin{flushright}
    \textit{Uso Docente:} \quad $\square$ Seguro \quad $\square$ Parcial \quad $\square$ Débil
\end{flushright}

\end{document}
